\chapter{Introduction, Motivation, and Thesis Scope}

Pakistan's agricultural sector faces a structurally difficult decision problem: managers must choose crops and input levels under uncertain weather, rising fertilizer prices, and delayed feedback from soil and yield outcomes. In practice, many farm-management decisions remain heuristic, calendar-driven, or based on short-horizon intuition rather than explicit long-term optimization. This creates a strong motivation for simulation-first decision support, especially where field experimentation is costly and risky. Official policy documents continue to place agriculture at the center of the national economy, while also emphasizing climate exposure, productivity pressure, and input-cost sensitivity \citep{pakistaneconomicsurvey2025}. In parallel, agricultural reinforcement learning (RL) has matured from a niche idea into a credible research direction for sequential decision-making in crop management \citep{gautron2022review}.

The present thesis is motivated by a gap between that global research direction and the concrete state of local, Pakistan-grounded tooling. Existing benchmark environments such as CropGym and gym-DSSAT have advanced the field, but the repository used in this thesis was originally based on CyclesGym because CyclesGym is more naturally aligned with multi-year crop-management questions, soil carry-over effects, and rotation-aware reasoning \citep{kallenberg2023cropgym,gautron2022gymdssat,turchetta2022cyclesgym}. However, the original base repository was not ready to serve as a Pakistan-focused thesis artifact without substantial additional engineering. It needed localized price series, clearer scope boundaries, reporting outputs suitable for academic defense, and explicit traceability from code changes to thesis claims.

\section{Why a Simulation-First RL Thesis is Defensible}

Simulation-first RL is defensible in agriculture for two reasons. First, agricultural decision quality is usually evaluated only after a season or multiple seasons, which makes exhaustive field experimentation impractical. Second, crop growth models create a repeatable environment in which policy learning can be studied before any pilot deployment. This does not eliminate the need for field validation, but it changes the stage at which expensive field testing is required. Instead of using the field to discover whether a policy family is viable at all, simulation can be used to narrow the search space to a smaller number of plausible policies \citep{turchetta2022cyclesgym,gautron2022review}.

The thesis therefore adopts a research posture that is realistic rather than promotional. It does not claim that simulation replaces agronomy or field trials. Instead, it treats the simulator as a controlled environment for learning decision policies, benchmarking algorithmic configurations, and stress-testing reward designs. This distinction is especially important because the repository is not a generic agricultural AI platform. It is a specific RL-enabled wrapper around the CYCLES simulator, and the defensibility of the thesis depends on keeping every claim anchored to what that stack actually does.

\section{Pakistan Context and Input-Cost Motivation}

A key thesis argument is that economic localization matters. It is not enough to optimize yield using generic or imported assumptions if the deployed environment is intended to represent Pakistan-relevant decisions. The repository already contains a Pakistan weather file, Pakistan soil profile, and a generated yearly price series derived from official data channels including FAOSTAT and NFDC-linked fertilizer prices. Figures~\ref{fig:pk-crop-prices}, \ref{fig:pk-nutrient-prices}, and \ref{fig:pk-weather-coverage} summarize the localized signals that now shape the active thesis stack.

\begin{figure}[H]
    \centering
    \includegraphics[width=0.92\textwidth]{pakistan_crop_price_trends.pdf}
    \caption{Pakistan crop price series used by the thesis stack, generated from the repository's official-source-derived yearly price file.}
    \label{fig:pk-crop-prices}
\end{figure}

\begin{figure}[H]
    \centering
    \includegraphics[width=0.92\textwidth]{pakistan_nutrient_price_trends.pdf}
    \caption{Pakistan nutrient price series used by the thesis stack. These series motivate a cost-aware, rather than yield-only, optimization objective.}
    \label{fig:pk-nutrient-prices}
\end{figure}

\begin{figure}[H]
    \centering
    \includegraphics[width=0.92\textwidth]{pakistan_weather_coverage.pdf}
    \caption{Coverage of the Pakistan weather file wired into training and evaluation. The thesis uses this file as the local weather backbone for the implemented environments.}
    \label{fig:pk-weather-coverage}
\end{figure}

These figures do not by themselves prove agronomic optimality. Their role is narrower and more important: they show that the repository now operates on localized inputs instead of remaining tied to generic defaults. That localization matters because the thesis is explicitly framed around Pakistan, and local economic conditions change both the meaning of reward design and the interpretation of policy quality.

\begin{figure}[H]
    \centering
    \includegraphics[width=0.92\textwidth]{pakistan_monthly_climatology.pdf}
    \caption{Monthly climatology implied by the Pakistan weather file used by the thesis stack. The figure clarifies the seasonal structure against which weekly and yearly decisions are learned.}
    \label{fig:pk-monthly-climatology}
\end{figure}

The climatology perspective is important because the thesis is not optimizing in a stationary tabular world. The simulated environment reflects a strong seasonal structure, with rainfall concentration and temperature shifts that determine planting windows, crop development, and the timing sensitivity of nutrient application. A weekly RL policy that ignores this seasonal rhythm would not be scientifically meaningful even if it achieved a numerically high reward inside a narrow benchmark configuration.

\section{Operational Decomposition of the Thesis Problem}

The thesis decomposes the agricultural decision problem into three layers. The first layer is \emph{data realism}: weather, soil, and price inputs must be localized enough that the reward signal carries Pakistan-relevant meaning. The second layer is \emph{control design}: the environment must expose actions that correspond to recognizable management decisions such as crop choice, planting timing, and nutrient application. The third layer is \emph{evidence discipline}: the training pipeline must produce outputs that allow the thesis to reconstruct what happened during learning and evaluation rather than only the final return.

This decomposition is one reason the current thesis is organized as a systems thesis as much as an algorithm thesis. A weak data layer would make cost-aware claims fragile. A weak control layer would make the learned policy difficult to interpret agronomically. A weak evidence layer would make the results chapter vulnerable to committee criticism because the thesis could not explain why a particular score occurred. The repository changes described in later chapters should therefore be read as targeted interventions on these three layers rather than as miscellaneous software improvements.

\section{Problem Statement}

The core problem addressed in this thesis is the following: how can a CYCLES-based RL environment be transformed into a Pakistan-adapted, cost-aware experimental platform that supports credible long-horizon crop-management research without overstating what has and has not been implemented? The question is not just whether an RL agent can achieve a high score. The harder question is whether the surrounding system, including data sources, action design, reward formulation, reporting outputs, and experiment runner, is rigorous enough to support a thesis that will be examined on both engineering and scientific grounds.

This problem has several dimensions. At the algorithmic level, the repository must support policies for sequential management decisions such as fertilization and crop planning. At the data level, it must replace legacy assumptions with Pakistan-relevant weather, soil, and price signals. At the reporting level, it must expose not only aggregate returns but also the intermediate decisions that matter to a thesis committee: which crop was selected, how much nutrient was applied, what the cost decomposition was, and whether a planting decision respected seasonal windows. At the academic-writing level, it must distinguish between implemented capability, planned experimentation, and future work.

\section{Research Questions}

The thesis is structured around four operational research questions.

\begin{enumerate}[leftmargin=1.2cm]
    \item How can the original CyclesGym base repository be localized for Pakistan-relevant crop-management experimentation while preserving backward compatibility where possible?
    \item What architectural and reporting extensions are required to move from a largely nitrogen-centric workflow to a broader NPK-aware and hierarchical thesis stack?
    \item How should the empirical chapter be structured when the latest broad NPK matrix is fully defined and ready to run, but not yet completed at the time of the current draft build?
    \item Which thesis claims can be defended today as implemented contributions, and which must be retained as partial work or future extensions?
\end{enumerate}

These questions are intentionally phrased to keep the thesis aligned with the actual repository state. The work is not presented as a completed farmer-advisory deployment, nor as a finalized field-validated agronomic product. It is presented as a thesis-grade research stack that has been materially extended and prepared for a broad experiment campaign.

\section{Evidence Hierarchy Used in This Document}

Not all evidence used in the thesis has the same epistemic role. Primary external literature establishes why simulation-based RL, crop-growth models, and cost-aware formulations are legitimate research directions. Official Pakistan and NUST sources establish the contextual and institutional grounding of the work. Repository artifacts establish what has actually been implemented. Experiment outputs, once the full campaign is run, will establish comparative performance within the defined scope.

This hierarchy matters because it prevents category mistakes. A proposal slide cannot substitute for a code implementation. A unit test cannot substitute for a completed experiment matrix. A historical benchmark from an older stack cannot substitute for the final latest-NPK evidence. The writing strategy of the thesis deliberately follows this hierarchy so that each claim is backed by the right kind of source.

\section{Contributions of the Thesis}

The first contribution is technical localization. The repository now includes a Pakistan-oriented weather and soil setup, crop-calendar enforcement hooks, and yearly price series that shift the reward design away from imported default assumptions. The second contribution is representational expansion. The fertilization pipeline is no longer restricted to a simple nitrogen-only framing; the codebase now supports NPK-aware logic and price decomposition. The third contribution is hierarchical integration. The system now supports a yearly crop-planning layer coupled with within-season fertilization decisions, which better reflects the multiscale nature of real farm management. The fourth contribution is reproducibility and reporting. Standardized summary JSON files, reporting callbacks, and a unified experiment runner make the results chapter reproducible rather than anecdotal.

A fifth contribution is methodological honesty. This thesis does not collapse planning documents, code changes, and completed results into a single unqualified narrative. Instead, it explicitly separates the synopsis promises from the implemented repository scope. That distinction is essential because the synopsis included ambitions such as irrigation control and rice-wheat localization that are not yet implemented in the active repo. The thesis therefore includes a synopsis closure matrix so that each promise is classified as implemented, partial, or future work.

Taken together, these contributions mean that the thesis should be evaluated on two axes. The first is whether the engineering stack is meaningfully better aligned with Pakistan-grounded, cost-aware agricultural RL than the inherited base repository. The second is whether the document preserves a disciplined boundary between implemented capability and deferred research ambition. The current draft argues that the answer is yes on both axes, even though the final broad empirical campaign is still pending.

\input{tables/generated/synopsis_closure_matrix.tex}

\section{Scope Boundaries}

The most important boundary is that the thesis methods and claims reflect the current working implementation, not the broadest aspiration in the synopsis. The implemented empirical scope covers Pakistan-adapted weather and soil integration, crop planning, NPK-aware fertilization, hierarchical RL scaffolding, and the latest unified experiment matrix. The active crop-planning and hierarchical setup remains tied to the current maize-soy configuration already wired into the repository. By contrast, irrigation as a learned action, rice-specific localization, rice-wheat experiments, and formal field validation remain out of scope for the implemented results.

This bounded scope is a strength rather than a weakness. It prevents the thesis from claiming more than the repository can support, and it creates a clear roadmap for post-thesis work. In a field where engineering complexity often obscures what is truly implemented, scope clarity is part of scientific rigor.

The maize-soy decision should be read in this light. It is not presented as the unique or final agronomic target for Pakistan. It is presented as the currently stable experimental configuration for which the code, tests, runner defaults, and reporting pipeline are all aligned. In other words, the thesis prioritizes internal coherence over premature expansion to more ambitious but less stable crop configurations.

\section{Organization of the Thesis}

Chapter~2 establishes the literature and methodological background, with particular attention to crop-growth models, RL, constrained optimization, and related agricultural environments. Chapter~3 explains how the original base repository was transformed into the current Pakistan-adapted thesis stack. Chapter~4 formalizes the methodological and experimental design, including state, action, reward, matrix structure, and hypotheses. Chapter~5 documents the implementation, data provenance, reporting pipeline, and reproducibility controls. Chapter~6 reports the current execution status of the latest NPK matrix and contextualizes it against historical benchmark evidence from the earlier audited runs. Chapter~7 discusses limitations, closes the loop on synopsis promises, and outlines future work.
